\section{Results}
\label{sec:results}

\subsection{H-C1: Doctest Feasibility Scan}

We present the three-phase filtering funnel results from the scan of 10,000 Python files.

\subsubsection{Primary Finding: 31$\times$ Gap}

\begin{table}[h]
\centering
\caption{Three-phase filtering funnel results (10,000 Python files).}
\label{tab:funnel}
\begin{tabular}{llll l}
\toprule
\textbf{Metric} & \textbf{Count} & \textbf{Rate} & \textbf{Threshold} & \textbf{Status} \\
\midrule
n\_sampled & 10,000 & 100\% & --- & --- \\
Phase A (pattern) & 310 & 3.1\% & --- & --- \\
Phase B (AST) & 204 & 2.0\% & --- & --- \\
\textbf{Phase C (executable)} & \textbf{10} & \textbf{0.1\%} & \textbf{3.0\%} & \textbf{PIVOT} \\
\bottomrule
\end{tabular}
\end{table}

Figure~\ref{fig:gate_metrics} shows the three-phase rates against the 3\% SHOULD\_WORK
threshold. The pattern rate (3.1\%) meets the threshold; the executable rate (0.1\%) falls
31$\times$ below it.

Figure~\ref{fig:prevalence} visualizes the stacked breakdown: 96.9\% of files have no
doctest patterns; 1.1\% have patterns but fail AST parsing; 1.9\% pass Phase B but fail
Phase C; 0.1\% pass all three phases.

\subsubsection{Token Pool Infeasibility}

\begin{table}[h]
\centering
\caption{Token pool vs.\ SFT budget target.}
\label{tab:token}
\begin{tabular}{ll}
\toprule
\textbf{Metric} & \textbf{Value} \\
\midrule
Estimated executable-doctest files in full corpus & 12,960 \\
Estimated token pool & \textbf{0.004M tokens} \\
SFT budget target & 500M tokens \\
Ratio (actual / target) & \textbf{0.0008\%} \\
\bottomrule
\end{tabular}
\end{table}

Figure~\ref{fig:token_pool} shows the 125,000$\times$ gap between actual token pool
and SFT budget target.

\subsubsection{Gate Result}

\begin{verbatim}
GATE CHECK: executable_rate=0.001
STATUS: PIVOT (<1%) -- 2-condition H-E1 design activated
\end{verbatim}

\subsubsection{Phase C Failure Analysis}

Figure~\ref{fig:error_dist} shows Phase C failure types across the 194 files that passed
Phase B. Import errors (\texttt{ImportError}, \texttt{ModuleNotFoundError}) are the dominant
failure mode, confirming the import isolation hypothesis.

\begin{table}[h]
\centering
\caption{Phase C failure type distribution (194 files).}
\label{tab:errors}
\begin{tabular}{ll}
\toprule
\textbf{Failure Type} & \textbf{Fraction} \\
\midrule
import\_error & Dominant ($\sim$95\%) \\
exec\_error & Minor \\
test\_failed & Rare \\
timeout & Rare \\
\bottomrule
\end{tabular}
\end{table}

\subsubsection{Pipeline Performance}

\begin{table}[h]
\centering
\caption{Pipeline performance summary.}
\label{tab:perf}
\begin{tabular}{ll}
\toprule
\textbf{Metric} & \textbf{Value} \\
\midrule
Total files scanned & 10,000 \\
Scan duration & \textbf{129.8 seconds} \\
Throughput & $\sim$77 files/second \\
Workers & 4 (ProcessPoolExecutor) \\
Unit tests & \textbf{28/28 pass} \\
\bottomrule
\end{tabular}
\end{table}

\subsection{H-E1: SFT Comparison (Pending)}

The SFT training experiment is designed and pending execution. H-E1 (Qwen2.5-Coder-1.5B,
compile-only vs.\ unfiltered, 500M tokens, HumanEval/MBPP pass@1) is provided here as a
fully-specified design for reproducibility and as designated future work; it is not a current
claim of this paper. The empirical outcome --- including whether compile-only filtering provides
any marginal benefit over existing heuristic filters --- will be determined by running the
experiment. Note that if The Stack's existing heuristic filters (line length, alphanumeric
fraction, encoding) already exclude the majority of syntactically invalid files, the marginal
improvement from adding \texttt{compile()} may be small.

\subsection{File Size Analysis}

Figure~\ref{fig:file_size} shows token count distributions for executable-doctest-bearing
files (n=10) versus non-executable files (n=9,990). Both distributions exhibit similar
profiles, ruling out file size as a confound.
